% trading-function.tex — Linearized power² trading function derivation
% Kontrol: prove_cfmm_tradingFunction*, prove_cfmm_coordinateTransform*

\section{Trading Function}\label{sec:trading-function}

\subsection{Reserve Derivation}

Starting from the LP payoff $V(p) = (\pu^2 - p^2)/\pl^2$ for $p \in [\pl, \pu]$,
we derive reserves via the Angeris replication theorem \cite{angeris2021replicating}.

\begin{definition}[Reserves]\label{def:reserves}
The risky (virtual concentration) and numeraire (USDC) reserves per unit liquidity:
\begin{align}
  x(p) &= -V'(p) = \frac{2p}{\pl^2} \label{eq:risky-reserve} \\
  y(p) &= V(p) + p \cdot x(p) = \frac{\pu^2 + p^2}{\pl^2} \label{eq:numeraire-reserve}
\end{align}
for $p \in [\pl, \pu]$, with boundary extensions:
$x = 2\pl/\pl^2$ and $y = (\pu^2 + \pl^2)/\pl^2$ for $p \leq \pl$;
$x = 2\pu/\pl^2$ and $y = 2\pu^2/\pl^2$ for $p \geq \pu$.
\end{definition}

\begin{proposition}[Reserve Verification]\label{prop:reserve-verify}
$V(p) = y(p) - p \cdot x(p)$ holds exactly:
\[
  \frac{\pu^2 + p^2}{\pl^2} - p \cdot \frac{2p}{\pl^2}
  = \frac{\pu^2 + p^2 - 2p^2}{\pl^2}
  = \frac{\pu^2 - p^2}{\pl^2}
  = V(p) \qquad\checkmark
\]
\end{proposition}

\subsection{Coordinate Transform: $u = x^2$}\label{subsec:coordinate-transform}

The natural trading function $y = \pu^2/\pl^2 + (\pl^2/4)x^2$ is quadratic in $x$.
This creates a 2-homogeneous invariant, incompatible with O(1) tick-based
liquidity aggregation (Uniswap V3-style $\Lact \mathrel{+}= \Lnet$ at tick crossings).

\begin{definition}[Linearization]\label{def:linearization}
Define the transformed risky reserve $u = x^2$. Then:
\begin{equation}\label{eq:linearized-invariant}
  \boxed{\psi(u, y) = y - \frac{\pl^2}{4}\, u}
\end{equation}
with invariant constant $C_1 = \pu^2/\pl^2$ per unit liquidity.
\end{definition}

\begin{theorem}[Linearity and 1-Homogeneity]\label{thm:linearity}
$\psi(u, y)$ is linear in $(u, y)$. With liquidity $L$:
\[
  u_{\text{total}} = L \cdot \tilde{u}(p), \quad
  y_{\text{total}} = L \cdot \tilde{y}(p)
\]
where $\tilde{u}(p) = 4p^2/\pl^4$ and $\tilde{y}(p) = (\pu^2 + p^2)/\pl^2$
are per-unit-liquidity reserves. The invariant scales as:
\[
  \psi(u_{\text{total}}, y_{\text{total}}) = L \cdot C_1
\]
which is 1-homogeneous in $L$.
\end{theorem}

\begin{proof}
$\psi(L\tilde{u}, L\tilde{y}) = L\tilde{y} - (\pl^2/4) \cdot L\tilde{u}
= L(\tilde{y} - (\pl^2/4)\tilde{u}) = L \cdot C_1$. \qedhere
\end{proof}

\begin{remark}[Concavity]
The Hessian of $\psi$ has eigenvalues $-\pl^2/2$ and $0$, both $\leq 0$.
Hence $\psi$ is concave, and the trading set
$S = \{(u, y) : \psi(u, y) \geq L C_1\}$ is convex
(Geometry paper \cite{angeris2023geometry}, Theorem~1).
\end{remark}

\subsection{Spot Price}

\begin{definition}[Spot Price from Reserves]\label{def:spot-price}
\begin{equation}\label{eq:spot-price}
  p = \frac{\pl^2}{2}\sqrt{\frac{u}{L}}
  = \frac{\pl^2}{2} \cdot \frac{x}{L}
  \quad\text{(in } x \text{ space: } p = \pl^2 x / (2L)\text{)}
\end{equation}
where the second equality uses $x = \sqrt{u}$.
\end{definition}

\begin{proof}[Verification]
From the trading function, $-\partial\psi/\partial u \big/ \partial\psi/\partial y
= -(-\pl^2/4) / 1 = \pl^2/4$. This is the marginal price in $(u, y)$ space.
In $(x, y)$ space, $du = 2x\, dx$, so the marginal price is
$dy/dx = (\pl^2/4) \cdot 2x = \pl^2 x/2 = p$, matching the reserve
derivation $x(p) = 2p/\pl^2 \Leftrightarrow p = \pl^2 x/2$. \qedhere
\end{proof}

\subsection{Swap Mechanics}

A swap in the insurance CFMM operates in $(u, y)$ space internally, with
$x \leftrightarrow u$ conversion at the boundary.

\begin{definition}[Swap]\label{def:swap}
Trader sells $\Delta x$ units of concentration exposure:
\begin{align}
  \Delta u &= (x + \Delta x)^2 - x^2 = 2x\,\Delta x + (\Delta x)^2 \label{eq:delta-u} \\
  \Delta y &= \frac{\pl^2}{4}\, \Delta u \label{eq:delta-y}
\end{align}
The effective price $\Delta y / \Delta x = (\pl^2/4)(2x + \Delta x)$ depends on
the current state $x$ and trade size $\Delta x$, providing natural slippage
despite the internal invariant being linear.
\end{definition}

\begin{remark}[O(1) Complexity]
Within a tick: the swap is a single multiplication (\cref{eq:delta-y}) plus
one square root for $x = \sqrt{u}$ conversion (\texttt{FixedPointMathLib.sqrt}).
Tick crossing: $\Lact \mathrel{+}= \Lnet$, a single addition.
Total per-swap cost: O(ticks crossed), same as Uniswap V3.
\end{remark}

\subsection{Liquidity Provisioning}

\begin{definition}[Single-Sided USDC Deposit]\label{def:single-sided}
An underwriter deposits $\Delta y$ USDC at current price $p$:
\begin{align}
  \Delta L &= \frac{\Delta y \cdot \pl^2}{\pu^2 + p^2} \label{eq:mint-L} \\
  \Delta u &= \Delta L \cdot \frac{4p^2}{\pl^4}
             = \frac{4p^2 \cdot \Delta y}{\pl^2(\pu^2 + p^2)} \label{eq:mint-u}
\end{align}
The virtual $\Delta u$ is computed, not deposited. Only USDC enters the contract.
\end{definition}

\begin{proposition}[Invariant Preservation on Mint]\label{prop:mint-invariant}
After minting: $\psi(u + \Delta u, y + \Delta y) = (L + \Delta L) \cdot C_1$.
\end{proposition}

\begin{proof}
$\psi(u + \Delta u, y + \Delta y) = (y + \Delta y) - (\pl^2/4)(u + \Delta u)
= \psi(u, y) + \Delta y - (\pl^2/4)\Delta u
= L C_1 + \Delta L \cdot \tilde{y} - (\pl^2/4)\Delta L \cdot \tilde{u}
= L C_1 + \Delta L \cdot C_1 = (L + \Delta L) C_1$. \qedhere
\end{proof}
